\documentclass[11pt,a4paper]{report}

\usepackage[top=0.5in, bottom=1in, left=0.8in, right=0.8in, headsep=0.in, centering]{geometry}
\usepackage{graphicx}
\usepackage{amsmath}
\usepackage{amsfonts}
\usepackage{amssymb}
\usepackage{amsthm}
\usepackage{bookmark}
\usepackage{enumitem}
\usepackage{float}

\newtheorem{theorem}{Theorem}[section]
\newtheorem{lemma}[theorem]{Lemma}
\newtheorem{proposition}[theorem]{Proposition}
\newtheorem{corollary}[theorem]{Corollary}
\theoremstyle{definition}
\newtheorem{definition}[theorem]{Definition}
\newtheorem{example}[theorem]{Example}
\newtheorem{remark}[theorem]{Remark}
\newtheorem{assumption}[theorem]{Assumption}
\newtheorem{problem}[theorem]{Problem}
\newtheorem{property}[theorem]{Property}

\def\Xint#1{\mathchoice
{\XXint\displaystyle\textstyle{#1}}%
{\XXint\textstyle\scriptstyle{#1}}%
{\XXint\scriptstyle\scriptscriptstyle{#1}}%
{\XXint\scriptscriptstyle\scriptscriptstyle{#1}}%
\!\int}
\def\XXint#1#2#3{{\setbox0=\hbox{$#1{#2#3}{\int}$}
\vcenter{\hbox{$#2#3$}}\kern-.5\wd0}}
\def\ddashint{\Xint=}
\def\dashint{\Xint-}

\begin{document}
\begin{titlepage}
	\centering
	{\scshape\LARGE PDE Notes \par}
	\vspace{1cm}
	{\scshape\Large \par}
	\vspace{1.5cm}
	{\huge\bfseries Notes on Inequalities and Embeddings\par}
	\vspace{2cm}
	{\Large\itshape Kiyuob Jung\par}
	\vspace{2cm}
	{\Large\bfseries (in progress)\par}
	\vspace{5cm}
	
	\begin{figure}[H] 
	\centering 
	\includegraphics[scale=0.4]{Illustration/Logo.png} 
	\end{figure}
	\vspace{2cm}

	Department of Mathematics,\par
	Kyungpook National University
	
	\vfill

% Bottom of the page
	{\large \today\par}
\end{titlepage}

\newpage

\tableofcontents

\newpage

\subsection*{Preface}

Todo.


\subsection*{Abbreviation}

Since these notes is not for formal research, we almost always employ the following abbreviations:

\noindent $\diamond$ s.t. : s.t.

\noindent $\diamond$ TFAE : The following are equivalent


\subsection*{Notation}

Let a set $X$. 
We employ the following notations:

\noindent $\diamond$ $\mathbb{R}^{n}$ : $n$-dimensional real Euclidean space, $\mathbb{R}=\mathbb{R}^{1}$

\noindent $\diamond$ $\mathbb{C}^{n}$ : $n$-dimensional complex space, $\mathbb{C}=\mathbb{C}^{1}$

\noindent $\diamond$ $\mathbb{K}$ : either $\mathbb{R}$ or $\mathbb{C}$.

\noindent $\diamond$ $\partial X$ : boundary of $X$.

\noindent $\diamond$ $\forall$ : for all.

\chapter{Definitions and notations} 

In this note, we always denote $x=(x_1, x_2, \ldots, x_n)$ to be a point in $\mathbb{R}^{n}$.

\section{Definitions}


Here, we collect definitions for averages of a function.

\begin{definition}
	Let $f \in L^1(U)$ with an open set $U \subset \mathbb{R}^{n}$.
	\begin{enumerate}[label=(\alph*)] 
	\rm\item An average of $f$ over set $E$ is 
	\begin{equation*} 
		\dashint_{E} f d x:=\frac{1}{\operatorname{meas}(E)} \int_{E} f(x) d x.
	\end{equation*}
	\rm\item An average of $f$ over the ball $B_r(x_0)$ is 
	\begin{equation*} 
		\dashint_{B_r(x_0)} f d x:= \frac{1}{\alpha(n) r^{n}} \int_{B_r(x_0)} f(x) d x.
	\end{equation*}
	\rm\item An average of $f$ over the sphere $\partial B_r(x_0)$ is 
	\begin{equation*} 
		\dashint_{\partial B_r(x_0)} f d S :=\frac{1}{n\alpha(n) r^{n-1}} \int_{\partial B_r(x_0)} f(x) d S.
	\end{equation*}
	\end{enumerate}
\end{definition}

As for integrability, TODO.

\begin{definition}
    Let $p>1$ and define $q \in \mathbb{R}$ by 
    \begin{equation*} 
        \frac{1}{p} + \frac{1}{q} = 1.
    \end{equation*}
    Then $p$ and $q$ are called \emph{conjugate exponents}.
\end{definition}

\begin{remark}
    A simple calculation shows the following:
    \begin{enumerate}[label=(\roman*)] 
    \rm\item $pq = p + q$,
    \rm\item $\displaystyle 1 = \frac{p+q}{pq}$,
    \rm\item $(p-1)(q-1)=1$,
    \rm\item $\displaystyle q= \frac{p}{p-1}$.
    \end{enumerate}
\end{remark}

\chapter{Basic properties}


\section{Change of variable}

\begin{proposition}
	Let $f:U \to \mathbb{R}$ with an open set $U \subset \mathbb{R}^{n}$.
	\begin{enumerate}[label=(\alph*)] 
	\rm\item $\displaystyle \int_{B_r(x_0)} f(ax + b) dx = \frac{1}{r^n}\int_{B_{ar}(x_0 + b)} f(x) dx$.
	\rm\item $\displaystyle \dashint_{B_r(x_0)} f(ax + b) dx = \dashint_{B_{ar}(x_0 + b)} f(x) dx$.
	\end{enumerate}
\end{proposition}

\begin{proof} 
	TODO 

	To show (b), let $\tilde{x} := ax + b$. 
	Then $d \tilde{x} = r^n dx$.
	Also, since $\left\vert \tilde{x} - b \right\vert = \left\vert ax \right\vert < \left\vert a \right\vert r$, we get $\tilde{x} \in B_{\left\vert a \right\vert r}(x_0 + b)$.
	Hence, 
	\begin{align*}
		\dashint_{B_r(x_0)} f(ax + b) dx &= \frac{1}{\alpha(n) r^{n}} \int_{B_r(x_0)} f(x) d x \\
		&= \dashint_{B_{ar}(x_0 + b)} f(\tilde{x}) d\tilde{x} \\
		&= \dashint_{B_{ar}(x_0 + b)} f(\tilde{x}) d\tilde{x}.
	\end{align*}
\end{proof}

\begin{remark}
	Convolution. TODO
\end{remark}

\section{Coordinates}

\subsection{Polar coordinates}

\begin{proposition}

Let $f: \mathbb{R}^{n} \rightarrow \mathbb{R}$ be continuous and summable. 
Then
\begin{enumerate}[label=(\alph*)]
\rm\item $\displaystyle \int_{\mathbb{R}^{n}} f(x) d x=\int_{0}^{\infty}\left(\int_{\partial B_r(x_{0})} f(x) d\mathcal{S}\right) d r \quad \forall x_{0} \in \mathbb{R}^{n}$.
\rm\item $\displaystyle \frac{d}{d r}\left(\int_{B_r(x_{0})} f(x) d x\right)=\int_{\partial B_r(x_{0})} f(x) d\mathcal{S} \quad \forall r>0$.
\end{enumerate}
\end{proposition}

\begin{proposition}

	Let $f: \mathbb{R}^{n} \rightarrow \mathbb{R}$ be continuous and summable. 
	Then
	\begin{enumerate}[label=(\alph*)]
	\rm\item $\displaystyle \int_{B_{\varepsilon}(0)} f(x) dx = \int_{0}^{\varepsilon} \left( \int_{\partial B_r(0)} f(x) d\mathcal{S}(x) \right) dr \quad \forall \varepsilon>0$.
	\end{enumerate}
	\end{proposition}

\chapter{Inequalities} 

\section{Power inequalities}

\begin{theorem}
	The following statements holds.
	\begin{enumerate}[label=(\alph*)] 
		\rm\item $1 + x \leq e^{x} \quad \forall x \in \mathbb{R}$.
		\rm\item (Cauchy's inequality) \\[0.1cm]
		$\displaystyle xy \leq \frac{x^{2}}{2}+\frac{y^{2}}{2} \quad \forall x, y \in \mathbb{R}$ 
		\rm\item $\displaystyle e^{(x + y)/2} < \frac{e^{y} - e^{x}}{y - x} \quad \forall x, y \in \mathbb{R}$ with $x \neq y$.
	\end{enumerate}
	
\end{theorem}

\begin{theorem}
	The following statements holds.
	\begin{enumerate}[label=(\alph*)] 
		\rm\item $\displaystyle \left(\frac{1}{e}\right)^{\frac{1}{e}} \leq x^{x} \quad \forall x > 0$.
		\rm\item $\displaystyle x \leq x^{x^{x}} \quad \forall x > 0$.
		\rm\item $1 < x^{y}+y^{x}  \quad \forall x, y > 0$
		\rm\item $x^{y}+y^{x} \leq x^{x}+y^{y} \quad \forall x, y > 0$
		\rm\item $x^{e y}+y^{e x} \leq x^{e x}+y^{e y} \quad \forall x, y > 0$
		\rm\item $\displaystyle \frac{1-\frac{1}{x^{y}}}{y} \leq \ln (x) \leq \frac{x^{y} - 1}{y} \quad \forall x, y > 0$. The upper and lower bounds converge to $\ln (x)$ as $y \rightarrow 0$.
		\rm\item $2 < (x+y)^{z}+(x+z)^{y}+(y+z)^{x} \quad \forall x, y, z > 0$.
		\rm\item $x^{2 y}+y^{2 z}+z^{2 x} \leq x^{2 x}+y^{2 y}+z^{2 z} \quad \forall x, y, z > 0$.
		\rm\item $(x y z)^{(x+y+z) / 3} \leq x^{x} y^{y} z^{z} \quad \forall x, y, z > 0$
	\end{enumerate}
	
\end{theorem}

\begin{theorem}
	The following statements holds.
	\begin{enumerate}[label=(\alph*)] 
		\rm\item (Cauchy's inequality with $\varepsilon$)\\[0.1cm]
		$\displaystyle xy \leq \epsilon x^{2}+\frac{y^{2}}{4 \varepsilon} \quad \forall x, y>0$ , $\forall \varepsilon>0$.
	\end{enumerate}
\end{theorem}

\begin{theorem}
	The following statements holds.
	\begin{enumerate}[label=(\alph*)] 
		\rm\item $(x+y)^{p} < x^{p}+y^{p} \quad \forall x, y > 0$, $\forall p \in (0, 1)$.
		\rm\item $(x + y)^{p} \leq 2^{p - 1}\left(x^{p} + y^{p}\right) \quad \forall x, y \geq 0$, $\forall p \in [1, \infty)$.
	\end{enumerate}
\end{theorem}


\chapter{Embeddings}


\section{Sobolev Embedding}

In this section, we deal with embeddings of diverse Sobolev spaces into others.
Given a Sobolev space, it automatically belongs to certain other space, depending on the relationship between the integrability $p$ and the dimension $n$.\footnote{The regularity of domains also affects the inclusion.}
There are three cases:
\begin{align*}
    &p \in [1, n), \\
    &p = n, \\
    &p \in (n, \infty].
\end{align*}
In particular, the second case $p = n$ is called the \emph{borderline case}. 
What are we trying to obtain from the Sobolev embedding theory?
Broadly speaking, given a Sobolev space $W^{k, p}$, imbeddings of $W^{k, p}$ target two types of Banach spaces: either another Sobolev space $W^{j, q}$ or H$\ddot{\rm o}$lder spaces $C^{j, \alpha}$ for some constants $j \leq k$, $q \geq p$, and $0 \leq \alpha \leq 1$.

\subsection{The case $1 \leq p < n$}

\begin{definition}
    If $1 \leq p < n$, the \emph{Sobolev conjugate} $p^{\ast}$ of $p$ is defined by 
    \begin{equation*} 
        p^{\ast} := \frac{np}{n-p}.
    \end{equation*}
\end{definition}

\begin{remark}
    A simple calculation shows the following:
    \begin{enumerate}[label=(\roman*)] 
    \rm\item $p^{\ast} > p$,
    \rm\item $\displaystyle \frac{1}{p^{\ast}} = \frac{1}{p} - \frac{1}{n}$,
    \rm\item $p^{\ast} \to \infty$ as $p \to n$.
    \end{enumerate}
\end{remark}

\begin{theorem}
    \text{\rm (Gagliardo-Nirenberg-Sobolev inequality)}
    If $1 \leq p < n$, then there exists a constant $C$, depending only on $p$ and $n$, such that
    \begin{align*}
    \|u\|_{L^{p^*}\left(\mathbb{R}^n\right)} \leq C\|D u\|_{L^p\left(\mathbb{R}^n\right)},
    \end{align*}
    for all $u \in C_c^1\left(\mathbb{R}^n\right)$.
\end{theorem}

\begin{theorem}  
    \text{\rm (Estimates for $W^{1, p}, 1 \leq p < n$)}
    Let $U$ be a bounded, open subset of $\mathbb{R}^n$, and suppose $\partial U$ is $C^1$. 
    Assume $1 \leq p < n$ and $u \in W^{1, p}(U)$. 
    Then $u \in L^{p^*}(U)$, with the estimate
    \begin{align*}
    \|u\|_{L^{p^*}(U)} \leq C\|u\|_{W^{1, p}(U)},
    \end{align*}
    the constant $C$ depending only on $p, n$, and $U$.
\end{theorem}

\subsection{The case $p=n$}

\subsection{The case $n < p \leq \infty$}

For convenience as in the notation for the Sobolev conjugate, we write 
\begin{equation*} 
    \gamma:=1- \frac{n}{p},
\end{equation*}
whenever $n < p \leq \infty$.

\begin{theorem}
    \text{\rm (Morrey's inequality)}
    If $n < p \leq \infty$, then there exists a constant $C$, depending only on $p$ and $n$, such that
    \begin{align*}
    \|u\|_{C^{0, \gamma}\left(\mathbb{R}^n\right)} \leq C\|u\|_{W^{1, p}\left(\mathbb{R}^n\right)}
    \end{align*}
    for all $u \in C^1\left(\mathbb{R}^n\right)$.
\end{theorem}

\begin{theorem}
    \text{\rm (Estimates for $W^{1, p}$, $n<p \leq \infty$)}
    Let $U$ be a bounded, open subset of $\mathbb{R}^n$, and suppose $\partial U$ is $C^1$. Assume $n < p \leq \infty$ and $u \in W^{1, p}(U)$. 
    Then $u$ has a version $u^{\ast} \in C^{0, \gamma}(\overline{U})$ with with the estimate
    \begin{align*}
        \left\|u^*\right\|_{C^{0, \gamma}(\bar{U})} \leq C\|u\|_{W^{1, p}(U)},
    \end{align*}
    the constant $C$ depending only on $p, n$ and $U$.
\end{theorem}

\subsection{The case $p = \infty$}

\begin{theorem}
    \text{\rm (Characterization of $W^{1, \infty}$)}
    Let $U$ be open and bounded, with $\partial U$ of class $C^1$. 
    Then $u: U \rightarrow \mathbb{R}$ is Lipschitz continuous if and only if $u \in W^{1, \infty}(U)$
\end{theorem}

\end{document}